\section*{Training 1 · 30 points}
Write $P(p)=p_1^{1/4}p_2^{3/4}$. This is a valid expenditure function, for
example for the normalised Cobb--Douglas index
\[
 u(x)=\frac4{3^{3/4}}x_1^{1/4}x_2^{3/4}.
\]
Students need not reconstruct that primitive index.

\paragraph{(a), 6 points.}
\[
 h_1=\frac{\bar uP}{4p_1},\qquad
 h_2=\frac{3\bar uP}{4p_2},\qquad
 p_1h_1+p_2h_2=\bar uP=e.
\]
Suggested split: 2 points per quantity, 2 for the spending identity.

\paragraph{(b), 8 points.}
Solving $m=e(p,\bar u)$ gives $v=m/P$. Since
\[
 v_m=\frac1P,\quad v_{p_1}=-\frac{m}{4p_1P},\quad
 v_{p_2}=-\frac{3m}{4p_2P},
\]
Roy yields $x_1=m/(4p_1)$ and $x_2=3m/(4p_2)$. Substituting
$\bar u=m/P$ into h gives exactly those quantities. Suggested split:
2 for v, 4 for Roy and the demands, 2 for the parameter match.

\paragraph{(c), 10 points.}
Initially $P^0=1$, $x^0=(4,12)$ and $u^0=16$. At the new prices
$P^1=2$, $x^1=(1/4,12)$ and $u^1=8$.
With positive-loss conventions,
\[
 CV=e(p^1,u^0)-16=32-16=16,
\]
\[
 EV=16-e(p^0,u^1)=16-8=8,
\]
\[
 L_{CS}=\int_1^{16}\frac4s\,ds=4\log16=16\log2\approx11.0904.
\]
Thus $EV<L_{CS}<CV$. Both goods are normal. At the compensated new
income 32, the bundle is $(1/2,24)$, which restores utility 16.
Suggested split: 2 for bundles/utilities, 3 CV, 3 EV, 2 CS and comparison.

\paragraph{(d), 6 points.}
The old basket costs $16\cdot4+12=76$ at new prices. Making it affordable
adds $76-16=60$, whereas CV adds only 16. With income 76 and the new
prices, indirect utility is $76/2=38>16$; reoptimisation selects
$(19/16,57)$. Keeping a basket affordable generally overcompensates compared
with the cheapest way to restore utility. Suggested split: 2 for the
payment, 2 for the CV comparison, 2 for new utility and its interpretation.
